\chapter{Power Markets I: Design}\label{ch:m3:power-markets-design}

On Sunday 11 May 2025, between 13:00 and 14:00, German solar panels produced 44
gigawatts while the whole country consumed 40. The day-ahead auction had cleared that
hour the day before at minus 250 euros a megawatt-hour: producers paid to deliver
electricity, and consumers were paid to take it. Five months earlier, at 17:00 on a
windless 12 December 2024, the same auction had cleared at 936 euros. Electricity is the one
commodity that must be produced at the moment it is consumed, and power markets are
designed around that fact. This chapter describes the design: the \omterm{def:m3:power-markets-design:mtu}{market time unit},
the day-ahead auction and how its algorithm clears, the \omterm{def:m3:power-markets-design:merit}{merit order} that sets the
price, zones and nodes, and the markets for capacity and reserves that energy prices
alone do not pay for.

\section{Electricity as a commodity}

Electricity can hardly be stored at scale, and supply and demand on a grid must match
at every instant or its frequency drifts from its nominal 50 or 60 hertz. Markets
therefore trade energy for delivery over short, fixed intervals, and the grid operator
balances whatever remains in real time.

\begin{definition}[Market time unit]\label{def:m3:power-markets-design:mtu}
The \emph{market time unit}\index{market time unit} (MTU) is the delivery interval of
a power market's products: an energy product is a constant power delivered over one
MTU, so a quantity of energy in MWh per MTU.
\end{definition}

A power market is a sequence of markets for the same MTU: forwards and futures for
months and years ahead, the day-ahead auction, continuous intraday trading
(\cref{ch:m3:power-markets-trading}), and the operator's balancing in real time. The
day-ahead auction is the anchor: most volume in Europe is priced there, and most
forwards settle against its prices.

\section{The day-ahead auction and its clearing algorithm}

\begin{definition}[Day-ahead market]\label{def:m3:power-markets-design:da}
The \emph{day-ahead market}\index{day-ahead market} is a daily sealed-bid call auction,
held the day before delivery, that sets for each MTU of the next day one
uniform-price clearing (One Quant Book 2, chapter 4) per \omterm{def:m3:power-markets-design:zone}{bidding zone} from the
buy and sell orders submitted before its gate closure.
\end{definition}

For each MTU, buyers submit a stepwise demand curve (how much they will buy at each
price) and sellers a stepwise supply curve. The algorithm chooses the accepted orders
that maximise welfare, the value of the accepted bids minus the cost of the accepted
offers, subject to the network constraints between zones, and sets in each zone the
price at which supply and demand cross. Every accepted order is paid or pays that
price; nobody pays their bid.

\begin{proposition}[The clearing price supports the allocation]\label{prop:m3:power-markets-design:support}
At the clearing price $p$ of a single zone with stepwise curves, every accepted bid has
a limit $\ge p$ and every accepted offer a limit $\le p$; no rejected bid has a limit
above $p$ and no rejected offer a limit below $p$. The accepted volume maximises welfare.
\end{proposition}

\begin{proof}
Serving bids in decreasing and offers in increasing price order and stopping when the
next bid is worth less than the next offer adds each unit of positive surplus and no
unit of negative surplus: that is the welfare maximum. The marginal order, partly
filled, sets $p$, which by construction lies between the limits of the accepted and the
rejected orders on each side.
\end{proof}

Real auctions add orders that span several MTUs.

\begin{definition}[Block order]\label{def:m3:power-markets-design:block}
A \emph{block order}\index{block order} is an order to buy or sell a quantity in each of
several MTUs, all or nothing, with a limit on the average price over those MTUs.
\end{definition}

Blocks let a thermal plant, with its start-up cost and minimum run time, bid its true
economics. They make clearing a combinatorial problem: a block may be rejected although
the prices it would have received average above its limit (a paradoxically rejected
block), but never accepted when they average below it. Europe's day-ahead auctions are
cleared by one algorithm across the whole continent.

\begin{definition}[Market coupling]\label{def:m3:power-markets-design:coupling}
\emph{Market coupling}\index{market coupling} is the joint clearing of the day-ahead
auctions of several zones by one algorithm, which allocates the cross-border
transmission capacity implicitly: electricity flows from the low-price zone to the
high-price zone until prices converge or the capacity is used up.
\end{definition}

\begin{dated}{2026-09}{European day-ahead coupling}\label{dat:m3:power-markets-design:sdac}
The Single Day-Ahead Coupling runs one auction at 12:00 CET for delivery the next day,
cleared by the EUPHEMIA algorithm across all participating \omterm{def:m3:power-markets-design:zone}{bidding zones}. It moved from
60-minute to 15-minute \omterm{def:m3:power-markets-design:mtu}{market time units} on trading day 30 September 2025, for delivery
from 1 October 2025. The harmonised clearing-price limits were $-\qty{500}{EUR/MWh}$ and
$+\qty{4000}{EUR/MWh}$; ACER's Decision 02/2026 revised the methodology, and the minimum
was lowered to $-\qty{600}{EUR/MWh}$ from trading day 28 May 2026.
\end{dated}

\section{The merit order, spark and dark spreads}

\begin{definition}[Merit order]\label{def:m3:power-markets-design:merit}
The \emph{merit order}\index{merit order} is the ranking of generating capacity by
short-run marginal cost, lowest first; with offers at marginal cost, the supply curve of
the auction is the merit order and the price is the marginal cost of the last plant
needed.
\end{definition}

A thermal plant's marginal cost is its fuel cost per MWh of electricity plus the cost of
the allowances for its emissions (\cref{ch:m3:emissions-and-environmental-markets}):
$(P_{\mathrm{fuel}} + e\,P_{\mathrm{CO_2}})/\eta$, with $\eta$ its efficiency and $e$ the
emissions per MWh of fuel. Wind and solar have almost none.

\begin{definition}[Spark spread, dark spread]\label{def:m3:power-markets-design:spark}
The \emph{spark spread}\index{spark spread} of a gas-fired plant of efficiency $\eta$ is
the power price minus the price of the gas it burns per MWh of electricity,
$P_{\mathrm{power}} - P_{\mathrm{gas}}/\eta$; the \emph{dark spread}\index{dark spread} is
the same for a coal-fired plant.
\end{definition}

\begin{omfigure}
\begin{tikzpicture}
\begin{axis}[width=11cm,height=5.2cm,
  xlabel={cumulative capacity, GW},ylabel={marginal cost, EUR/MWh},
  tick label style={font=\small},label style={font=\small},
  xmin=0,xmax=80,ymin=0,ymax=290,grid=major,grid style={gray!25},
  table/col sep=comma]
\addplot[omDef,very thick,const plot] table[x=gw,y=mc]{figdata/markets-3/05-power-markets-design/merit.csv};
\draw[omThm,thick,dashed] (axis cs:45,0) -- (axis cs:45,255);
\draw[omThm,thick,dashed] (axis cs:60,0) -- (axis cs:60,255);
\node[omThm,font=\footnotesize,anchor=east,align=right] at (axis cs:44.2,218) {demand\\45 GW};
\node[omThm,font=\footnotesize,anchor=west,align=left] at (axis cs:60.8,218) {demand\\60 GW};
\node[font=\footnotesize,anchor=south west] at (axis cs:0.3,4) {wind, solar};
\node[font=\footnotesize,anchor=south] at (axis cs:25.5,86) {lignite};
\node[font=\footnotesize,anchor=south] at (axis cs:38.5,93) {CCGT};
\node[font=\footnotesize,anchor=south] at (axis cs:56.5,93) {coal};
\node[font=\footnotesize,anchor=south] at (axis cs:69,132) {OCGT};
\node[font=\footnotesize,anchor=south] at (axis cs:74.5,252) {oil};
\end{axis}
\end{tikzpicture}

\omcaption{An illustrative \omterm{def:m3:power-markets-design:merit}{merit order} with 10~GW of wind and solar: gas at
\qty{35}{EUR/MWh} of fuel, coal at \qty{12}{}, lignite at \qty{5}{}, carbon at
\qty{70}{EUR/t}. At 45~GW of demand a combined-cycle gas plant sets the price
(\qty{89.35}{EUR/MWh}); at 60~GW a coal plant does (\qty{89.68}{}). Data: the chapter's
tutorial.}\label{fig:m3:power-markets-design:merit}
\end{omfigure}

With the fuels of \cref{fig:m3:power-markets-design:merit} and power at
\qty{100}{EUR/MWh}, a 55\%-efficient gas plant earns a \omterm{def:m3:power-markets-design:spark}{spark spread} of
\qty{36.36}{EUR/MWh} and a 40\%-efficient coal plant a \omterm{def:m3:power-markets-design:spark}{dark spread} of \qty{70.00}{} before
carbon: after carbon the coal plant is dearer (\cref{ch:m3:emissions-and-environmental-markets}).
Adding renewables shifts the whole stack to the right and lowers the price in every hour
in which they produce: with 30~GW of wind and solar and 60~GW of demand the price falls
from \qty{89.68}{} to \qty{89.35}{EUR/MWh}; with 55~GW, to \qty{10}{EUR/MWh}.

\section{Zonal and nodal pricing}

A network has limits, and a price must say where it applies.

\begin{definition}[Bidding zone]\label{def:m3:power-markets-design:zone}
A \emph{bidding zone}\index{bidding zone} is the largest geographic area within which
market participants can trade energy without allocating transmission capacity; the
\omterm{def:m3:power-markets-design:da}{day-ahead market} sets one price per zone and MTU.
\end{definition}

Europe prices by zones, mostly national; the United States' organised markets price by
nodes.

\begin{definition}[Locational marginal price, congestion rent]\label{def:m3:power-markets-design:lmp}
The \emph{locational marginal price}\index{locational marginal price} (LMP) at a node of a
network is the cost of serving one more MWh there at least cost, subject to the network's
limits: in US markets the sum of a system energy price, a congestion component and a loss
component. The \emph{congestion rent}\index{congestion rent} on a constrained line is the
price difference between its ends times the flow on it; it accrues to whoever holds the
line's capacity rights (in coupled markets, the grid operators).
\end{definition}

\begin{proposition}[Two zones, one line]\label{prop:m3:power-markets-design:twozones}
Two zones cleared jointly through a line of capacity $C$ have equal prices if the flow
that equalises them is at most $C$. Otherwise the flow is $C$ from the cheap zone to the
dear one, each zone clears alone with that flow as a price-taking export or import, and
the \omterm{def:m3:power-markets-design:lmp}{congestion rent} is $(p_B - p_A)\,C$.
\end{proposition}

\begin{proof}
Joint clearing maximises total welfare. If the unconstrained optimum's flow fits the line
it is feasible and optimal, with one price. If not, welfare is concave in the flow, so the
constrained optimum puts the flow at its bound; each zone then clears with the fixed flow,
and the prices differ by the shadow price of the line.
\end{proof}

\begin{omfigure}
\centering
\begin{tikzpicture}[font=\small,
  zone/.style={draw,thick,rounded corners=3pt,align=center,minimum width=3.6cm,minimum height=1.8cm},
  arr/.style={-{Stealth[length=2.5mm]},very thick}]
\node[zone,draw=omDef,fill=omDef!8] (a) at (0,0) {zone A: windy\\40 GW renewables\\demand 50 GW};
\node[zone,draw=omThm,fill=omThm!8] (b) at (8,0) {zone B: still\\5 GW renewables\\demand 55 GW};
\draw[arr,omProp] (a.east) -- node[above=3pt,font=\footnotesize] {line full: 3 GW} (b.west);
\node[omProp,font=\footnotesize,anchor=north,align=center] at (4,-0.3) {congestion rent:\\\qty{7.30}{EUR/MWh} $\times$ 3 GW};
\node[font=\footnotesize,anchor=north] at (0,-1.15) {price \qty{82.38}{EUR/MWh} (lignite)};
\node[font=\footnotesize,anchor=north] at (8,-1.15) {price \qty{89.68}{EUR/MWh} (coal)};
\node[font=\footnotesize,anchor=north] at (4,-1.85) {with a 20~GW line: one price, \qty{89.35}{EUR/MWh} (CCGT)};
\end{tikzpicture}

\omcaption{Two coupled zones with the \omterm{def:m3:power-markets-design:merit}{merit order} of
\cref{fig:m3:power-markets-design:merit}. A 3~GW line is congested and prices split; the
grid operators collect about \money{EUR}{21900} for the hour. With a 20~GW line prices converge. Data:
the chapter's tutorial.}\label{fig:m3:power-markets-design:zones}
\end{omfigure}

Zonal pricing hides congestion inside a zone, which the operator must then relieve by
redispatch, paying some plants to produce less and others more; nodal pricing reveals it
in the price but gives thousands of prices instead of one, and exposes every generator to
the local congestion of its node. The next chapter shows how traders hedge that exposure.

\section{Capacity markets, ancillary services and negative prices}

An energy-only market pays generators only for the MWh they sell. A plant that runs a few
hours a year must earn its fixed costs in those hours, at scarcity prices, which regulators
often cap.

\begin{definition}[Capacity market, ancillary service]\label{def:m3:power-markets-design:capacity}
A \emph{capacity market}\index{capacity market} pays resources for being available to
produce (or to reduce consumption) in future periods of system stress, through auctions
held years ahead, with penalties if they fail to deliver. An \emph{ancillary
service}\index{ancillary service} is a service other than energy that the grid operator
buys to keep the system secure: frequency reserves that respond in seconds to minutes,
voltage support, and the ability to restart the grid after a blackout.
\end{definition}

The price cap and the \omterm{def:m3:power-markets-design:capacity}{capacity market} are two answers to the same problem. The Texas
market, energy-only, set its price at its cap of \qty{9000}{\$/MWh} for days in February
2021 during a winter storm; the regulator's order was challenged in court and upheld by the
Supreme Court of Texas in 2024. The market monitor found that holding prices at the cap for 32
hours after load shedding ended, until the morning of 19 February, added USD~16 billion to the
market's costs.

At the other end, prices go below zero when supply that will not or cannot stop exceeds
demand plus export capacity. Three kinds of producers keep producing at negative prices:
plants whose shutdown and restart cost more than the loss of running through a few hours;
producers paid a subsidy per MWh produced, who produce as long as the price plus the
subsidy is positive; and plants that must run for other reasons (heat for a town,
system security). On 11 May 2025 German solar output exceeded the country's load in five
hours (\cref{fig:m3:power-markets-design:may11}). In the German data, 457 hours had a negative
average price in 2024 and 573 in 2025 (hourly averages, as the Bundesnetzagentur also counts
them; from October 2025 the auction clears quarter-hours), almost all in the middle of the day
(\cref{fig:m3:power-markets-design:neghours}).

\begin{omfigure}
\begin{tikzpicture}
\begin{axis}[width=11cm,height=4.8cm,axis y line*=left,
  xlabel={hour, German local time},ylabel={EUR/MWh},
  tick label style={font=\small},label style={font=\small},
  xmin=-0.5,xmax=23.5,ymin=-300,ymax=200,grid=major,grid style={gray!25},
  xtick={0,3,6,9,12,15,18,21},
  table/col sep=comma]
\addplot[omThm,very thick,const plot mark mid] table[x=hour,y=price]{figdata/markets-3/05-power-markets-design/may11.csv};\label{plot:m3:may11price}
\end{axis}
\begin{axis}[width=11cm,height=4.8cm,axis y line*=right,axis x line=none,
  ylabel={GW},ylabel near ticks,
  tick label style={font=\small},label style={font=\small},
  xmin=-0.5,xmax=23.5,ymin=0,ymax=75,
  legend columns=3,legend style={font=\footnotesize,at={(0.5,-0.33)},anchor=north,draw=none,fill=none,/tikz/every even column/.append style={column sep=8pt}},
  table/col sep=comma]
\addlegendimage{/pgfplots/refstyle=plot:m3:may11price}\addlegendentry{day-ahead price}
\addplot[omProp,thick] table[x=hour,y=solar]{figdata/markets-3/05-power-markets-design/may11.csv};\addlegendentry{solar output}
\addplot[omDef,thick,dashed] table[x=hour,y=load]{figdata/markets-3/05-power-markets-design/may11.csv};\addlegendentry{load}
\end{axis}
\end{tikzpicture}

\omcaption{Germany on Sunday 11 May 2025, hourly: the day-ahead price (left axis) fell to
$-\qty{250.32}{EUR/MWh}$ at 13:00 while solar output (right axis) exceeded load from 11:00
to 16:00. Data: Bundesnetzagentur | SMARD.de, CC BY 4.0.}\label{fig:m3:power-markets-design:may11}
\end{omfigure}

\begin{omfigure}
\begin{tikzpicture}
\begin{axis}[width=11cm,height=4.4cm,ybar,bar width=7pt,
  xlabel={hour, German local time},ylabel={negative hours},
  tick label style={font=\small},label style={font=\small},
  xmin=-0.7,xmax=23.7,ymin=0,ymax=190,grid=major,grid style={gray!25},
  xtick={0,3,6,9,12,15,18,21},
  table/col sep=comma]
\addplot[fill=omDef!60,draw=omDef] table[x=hour,y=count]{figdata/markets-3/05-power-markets-design/neghours.csv};
\end{axis}
\end{tikzpicture}

\omcaption{German day-ahead hours with a negative price, 2024 and 2025 together (1{,}030
hours), by local hour of delivery: the solar hours dominate. Data: Bundesnetzagentur |
SMARD.de, CC BY 4.0.}\label{fig:m3:power-markets-design:neghours}
\end{omfigure}

\section{Tutorial: clearing an hour and coupling two zones}\label{tut:m3:power-markets-design:clear}

\begin{tutorial}
\textbf{Goal.} Clear one MTU from stepwise orders, then couple two zones through a line.
\textbf{End state:} the prices of \cref{fig:m3:power-markets-design:merit,fig:m3:power-markets-design:zones}
and the negative hours of \cref{fig:m3:power-markets-design:neghours}.
\begin{tutsteps}
\item \textbf{One zone.} Merit-order matching and the marginal order's price.
\omcode{code/firm/dayahead/firm_dayahead.py}{30}{58}{Uniform-price clearing of one market time unit.}\label{lst:m3:power-markets-design:clear}
\item \textbf{Two zones.} \Cref{prop:m3:power-markets-design:twozones}.
\omcode{code/firm/dayahead/firm_dayahead.py}{61}{75}{Two zones coupled through one line.}\label{lst:m3:power-markets-design:couple}
\item \textbf{Run} \texttt{m3\_power.hour(60, 10)}, \texttt{two\_zones(3)},
\texttt{negative\_stats()} and \texttt{fig\_power.py}.
\end{tutsteps}
\textbf{What to change next.} Raise the carbon price until coal and gas swap places in the
\omterm{def:m3:power-markets-design:merit}{merit order}; add a block offer for a lignite plant over the solar hours and see whether it
is accepted.
\end{tutorial}

\section{Build: the day-ahead clearing engine}\label{bld:m3:power-markets-design:dayahead}

\begin{build}
\bfield{Purpose} The miniature firm bids plants and batteries into day-ahead auctions and
must predict the prices its own bids help set: it needs a clearing engine it can run on its
own forecasts of everyone else's curves.
\bfield{Interface} \texttt{Order(side, price, qty)}; \texttt{clear(orders) -> Result(price,
volume, welfare, fills)}; \texttt{couple(zone\_a, zone\_b, capacity)};
\texttt{Block(name, price, qty, mtus)}; \texttt{clear\_with\_blocks(hourly, blocks)};
\texttt{marginal\_cost}; \texttt{spark\_spread}; \texttt{dark\_spread}.
\bfield{Rules} Uniform price per zone and MTU; the partly filled marginal order sets the
price; curves that do not cross are an error; a block is all-or-nothing and never
paradoxically accepted; enumeration is limited to small block sets (the real algorithm
uses branch and bound).
\bfield{Acceptance tests} \texttt{code/firm/dayahead/tests/}: the marginal offer sets the
price; a negative price; congestion splits prices and earns the rent; convergence with a
large line; a dear block rejected.
\bfield{Stretch} More than two zones on a meshed network with flow-based constraints;
15-minute MTUs with hourly blocks; the price-limit rules of the dated box.
\end{build}

\begin{omsources}
\item EPEX SPOT and the Market Coupling Steering Committee, notices on the 15-minute MTU
in SDAC (2025); ENTSO-E, Single Day-ahead Coupling.
\item ACER Decision 02/2026 on the harmonised maximum and minimum clearing price for SDAC;
SDAC communication note of 7 May 2026.
\item PJM, \emph{Formation of Locational Marginal Pricing and its System Energy, Congestion
and Loss Components}.
\item Supreme Court of Texas, decision of June 2024 on the PUCT's February 2021 pricing
orders (via Reed Smith client note).
\item Potomac Economics (ERCOT's independent market monitor), recommendation to the PUCT,
Project No.\ 51812, 4 March 2021 (Internet Archive copy).
\item Bundesnetzagentur | SMARD.de: day-ahead prices DE-LU, generation and load, 2024--2025; Bundesnetzagentur, press release
on the electricity market in 2025.
\end{omsources}

\section{Exercises}

\begin{exercise}[$\star$]\label{exo:m3:power-markets-design:1}
Offers: 300 MWh at \qty{0}{}, 200 at \qty{40}{}, 200 at \qty{90}{}, 100 at \qty{150}{EUR/MWh}.
Demand bids 450 MWh at the price cap. Give the price and the accepted volume of each offer.
\end{exercise}

\begin{exercise}[$\star$]\label{exo:m3:power-markets-design:2}
A gas plant of efficiency 50\% burns gas at \qty{30}{EUR/MWh} and emits 0.202 t per MWh of
gas; carbon costs \qty{70}{EUR/t}. What is its marginal cost?
\end{exercise}

\begin{exercise}[$\star$]\label{exo:m3:power-markets-design:3}
Why does every accepted seller receive the clearing price rather than its own offer? What
would sellers do under pay-as-bid?
\end{exercise}

\begin{exercise}[$\star\star$]\label{exo:m3:power-markets-design:4}
In \cref{fig:m3:power-markets-design:zones}, what are the flow, the prices and the
\omterm{def:m3:power-markets-design:lmp}{congestion rent} if the line's capacity is 3~GW? What minimum capacity makes the prices
converge?
\end{exercise}

\begin{exercise}[$\star\star$]\label{exo:m3:power-markets-design:5}
A solar farm receives a premium of \qty{60}{EUR/MWh} on top of the market price for every
MWh it produces. Below what price does it stop? In which hours of 11 May 2025 would it
have stopped?
\end{exercise}

\begin{exercise}[$\star\star$]\label{exo:m3:power-markets-design:6}
Explain why nodal prices differ inside a zone that has one zonal price, and who pays for
the difference in a zonal system.
\end{exercise}

\begin{exercise}[$\star\star\star$]\label{exo:m3:power-markets-design:7}
\emph{Coding.} With \texttt{negative\_stats}, give the number of negative hours in 2024 and
in 2025, the local hour with the most, and the lowest price and its hour.
\end{exercise}

\begin{exercise}[$\star\star\star$]\label{exo:m3:power-markets-design:8}
\emph{Find the flaw.} ``Negative prices prove the market is broken: nobody should ever pay
to produce.''
\end{exercise}

\section{Problem: The Day with Too Much Sun}

\begin{problem}[{Weekend problem --- 11 May 2025 from a lignite unit}]\label{pb:m3:power-markets-design:1}
A lignite unit of the \omterm{def:m3:power-markets-design:merit}{merit order} of \cref{fig:m3:power-markets-design:merit} (marginal
cost \qty{82.38}{EUR/MWh}) can run no lower than 400~MW. On 11 May 2025 the German
day-ahead prices were negative from 09:00 to 18:00 local time.

\medskip\noindent\textbf{Part I --- The day.}
\begin{enumerate}
\item In how many hours was the price negative, and what was their average?
\item What was the lowest price, and when?
\item In which hours did solar output exceed load?
\item Where did the surplus go?
\item Why did the auction not clear at zero?
\end{enumerate}

\medskip\noindent\textbf{Part II --- The unit.}
\begin{enumerate}[resume]
\item What does the unit lose, relative to not producing, in one hour at minimum load at the
lowest price?
\item What does it lose over the nine negative hours at minimum load?
\item Why does it compare this loss with its restart cost?
\item State the rule for staying on.
\item What bid would express that rule in the auction?
\end{enumerate}

\medskip\noindent\textbf{Part III --- The subsidised producers.}
\begin{enumerate}[resume]
\item A solar farm with a premium of \qty{60}{EUR/MWh} produces in which negative hours?
\item What price does a large block of such farms put on the auction?
\item How have subsidy rules been changed to limit this?
\item Why does battery storage raise midday prices?
\item Which producers profit from negative prices?
\end{enumerate}

\medskip\noindent\textbf{Part IV --- Judgement.}
\begin{enumerate}[resume]
\item What would an interconnector twice as large have done to the German price that day?
\item Why are negative prices concentrated in spring?
\item What is the value of a flexible consumer on such a day?
\item State the \emph{named result}: the restart cost below which the unit should switch off.
\item In one sentence: what does a negative power price signal?
\end{enumerate}
\end{problem}

\section{Interview questions}

\begin{interviewq}[$\star$ \iqroles{trader}]\label{iq:m3:power-markets-design:1}
How is the day-ahead power price in a European zone set?
\end{interviewq}

\begin{interviewq}[$\star$ \iqroles{trader, researcher}]\label{iq:m3:power-markets-design:2}
What is the \omterm{def:m3:power-markets-design:merit}{merit order}, and how does adding wind change prices?
\end{interviewq}

\begin{interviewq}[$\star\star$ \iqroles{researcher}]\label{iq:m3:power-markets-design:3}
You want to forecast tomorrow's day-ahead price for each hour. What inputs matter, and
which is most important in spring?
\end{interviewq}

\begin{interviewq}[$\star\star$ \iqroles{trader, risk}]\label{iq:m3:power-markets-design:4}
Why can a power price be negative, and how deep can it go?
\end{interviewq}

\begin{interviewq}[$\star\star$ \iqroles{researcher}]\label{iq:m3:power-markets-design:5}
Compare zonal and nodal pricing from the point of view of a wind farm owner.
\end{interviewq}

\begin{interviewq}[$\star\star\star$ \iqroles{developer, researcher}]\label{iq:m3:power-markets-design:6}
Design a day-ahead clearing simulator for ten zones and 96 MTUs with \omterm{def:m3:power-markets-design:block}{block orders}. How do
you keep it fast enough to run thousands of scenarios?
\end{interviewq}
